\chapter{Venere e Marte Hanno lo Stesso Cervello}

\begin{quote}
	\textit{``La scienza non è né maschio né femmina. Ma gli scienziati, purtroppo, lo sono.''}\\
	\hfill Rita Levi-Montalcini
\end{quote}

Nel 2015, un gruppo di neuroscienziati dell'Università di Tel Aviv fece qualcosa che nessuno aveva tentato prima su quella scala: analizzò le risonanze magnetiche di oltre 1.400 cervelli umani, cercando di rispondere alla domanda più vecchia delle neuroscienze. Esiste davvero un ``cervello maschile'' e un ``cervello femminile''? La risposta, pubblicata sulla rivista \textit{Proceedings of the National Academy of Sciences}, sorprese molti: no. O meglio, non nel modo in cui la cultura popolare lo immagina.\footnote{Joel, D., et al. (2015). ``Sex beyond the genitalia: The human brain mosaic''. \textit{Proceedings of the National Academy of Sciences}, 112(50), 15468--15473.}

Quello che Daphna Joel e il suo team scoprirono è che ogni cervello è un mosaico unico. Esistono, sì, alcune strutture che in media differiscono tra i sessi (il volume di certe aree, la densità di alcune connessioni), ma queste differenze si distribuiscono come l'altezza: i gruppi si sovrappongono enormemente, e nessun cervello singolo è ``interamente maschile'' o ``interamente femminile''. È come se la natura avesse a disposizione una tavolozza di colori e, per ogni cervello, scegliesse a caso tra sfumature tradizionalmente associate all'uno o all'altro sesso, creando ogni volta un quadro irripetibile.

Perché questo è importante per un libro sui bias cognitivi? Perché la credenza in cervelli radicalmente diversi è essa stessa uno dei bias più potenti della nostra cultura. E perché le differenze reali, quelle piccole, misurabili, statistiche, raccontano una storia molto più affascinante di qualsiasi semplificazione da best-seller.

\section{La fiducia e il suo opposto}

Nel 2001, due economisti della University of California, Brad Barber e Terrance Odean, pubblicarono uno studio che avrebbe fatto storia nella finanza comportamentale. Analizzarono oltre 35.000 conti di investimento su un arco di sei anni e scoprirono un dato curioso: gli uomini facevano operazioni di compravendita il 45\% più frequentemente delle donne. Il risultato? I loro portafogli rendevano, in media, l'1,5\% in meno all'anno. Non poco, su un arco di decenni.\footnote{Barber, B. M., \& Odean, T. (2001). ``Boys will be boys: Gender, overconfidence, and common stock investment''. \textit{The Quarterly Journal of Economics}, 116(1), 261--292.}

Il meccanismo è quello che gli psicologi chiamano \textit{overconfidence bias}, eccesso di fiducia nelle proprie capacità. Ed è qui che emerge una delle poche differenze cognitive di genere che la ricerca ha documentato con una certa consistenza: gli uomini, in media, tendono a sovrastimare le proprie competenze; le donne, in media, tendono a sottostimarle.\footnote{Lundeberg, M. A., Fox, P. W., \& Puncochar, J. (1994). ``Highly confident but wrong: Gender differences and similarities in confidence judgments''. \textit{Journal of Educational Psychology}, 86(1), 114--121.}

Attenzione: ``in media'' è la parola chiave. Non stiamo parlando di leggi universali, ma di distribuzioni statistiche che si sovrappongono ampiamente. Esistono uomini pieni di dubbi e donne traboccanti di sicurezza. Ma quando si guardano i grandi numeri, la tendenza emerge con una certa regolarità, attraverso culture e contesti diversi.

Da dove viene questa asimmetria? La risposta, come quasi sempre in psicologia, è: da un intreccio di biologia e cultura così stretto che separare i fili è quasi impossibile. La psicologia evolutiva suggerisce che, nelle condizioni ancestrali, i maschi che esibivano sicurezza (anche se esagerata) avevano un vantaggio competitivo nell'accesso alle risorse e ai partner. Le femmine che calibravano il rischio con maggiore cautela proteggevano meglio la prole. Ma questa è una narrazione plausibile, non una certezza. La socializzazione conta almeno altrettanto: fin dalla prima infanzia, maschi e femmine ricevono messaggi diversi su quanto sia accettabile essere sicuri di sé, rischiare, sbagliare.\footnote{Eliot, L. (2009). \textit{Pink Brain, Blue Brain: How Small Differences Grow into Troublesome Gaps}. Houghton Mifflin Harcourt.}

Quello che possiamo dire con certezza è che entrambe le tendenze sono errori. L'eccesso di fiducia fa prendere decisioni avventate. Il difetto di fiducia fa perdere opportunità. Sono, per così dire, due modi diversi di sbagliare la mira: uno spara troppo in alto, l'altro troppo in basso. E nessuno dei due colpisce il bersaglio.

\section{La soglia invisibile della candidatura}

Nel 2014, un dato proveniente dalle risorse umane di Hewlett-Packard cominciò a circolare nel dibattito pubblico e non ha più smesso: secondo un'analisi interna, le donne tendevano a candidarsi per una promozione solo quando soddisfacevano il 100\% dei requisiti elencati. Gli uomini si candidavano già al 60\%.\footnote{Mohr, T. S. (2014). ``Why Women Don't Apply for Jobs Unless They're 100\% Qualified''. \textit{Harvard Business Review}.}

Il dato è stato talvolta semplificato oltre misura, ma il fenomeno sottostante è stato confermato da numerose ricerche successive. E il meccanismo è più sottile di quanto sembri a prima vista. Non si tratta semplicemente di donne ``timide'' e uomini ``coraggiosi''. Interviste approfondite hanno rivelato che le donne non si candidano perché interpretano i requisiti come criteri rigidi (``se non li soddisfo tutti, non sono qualificata''), mentre gli uomini li leggono come linee guida indicative (``ci provo, al massimo dicono no'').

La differenza non sta tanto nella fiducia in sé, quanto nella percezione delle regole. E qui entra in gioco un fattore che trascende la psicologia individuale: il contesto. Ricerche della psicologa sociale Brenda Major hanno mostrato che quando le donne operano in ambienti dove il feedback è chiaro e i criteri sono trasparenti, il gap di fiducia si riduce drasticamente. Non è il cervello femminile a essere ``difettoso'': è che il cervello femminile, in media, è più sensibile ai segnali sociali, e in ambienti ambigui o ostili, quei segnali dicono ``stai attenta''.\footnote{Major, B. (1989). ``Gender differences in comparisons and entitlement: Implications for comparable worth''. \textit{Journal of Social Issues}, 45(4), 99--115.}

C'è poi un aspetto che rende il quadro ancora più complesso: il doppio vincolo. La psicologa Alice Eagly lo ha documentato in decenni di ricerche sulla leadership.\footnote{Eagly, A. H., \& Karau, S. J. (2002). ``Role congruity theory of prejudice toward female leaders''. \textit{Psychological Review}, 109(3), 573--598.} Quando le donne adottano comportamenti assertivi, tipicamente premiati negli uomini (negoziare con fermezza, esprimere sicurezza, prendere decisioni rapide), vengono spesso penalizzate socialmente. ``Troppo aggressiva.'' ``Poco collaborativa.'' Gli stessi comportamenti che in un uomo verrebbero etichettati come ``leadership'' e ``determinazione''. È un vicolo cieco cognitivo: sei troppo timida se non osi, troppo dura se osi.

Questo non significa che la sottostima delle proprie capacità sia solo un prodotto culturale. Ma significa che non è possibile capirla guardando solo dentro il cranio. Bisogna guardare anche fuori, all'ambiente che quel cranio abita.

\section{Il mosaico della propensione al rischio}

Per decenni, una delle ``verità'' più consolidate delle scienze comportamentali è stata: gli uomini rischiano di più, le donne rischiano di meno. La realtà, come hanno mostrato ricerche più recenti, è considerevolmente più sfumata.

È vero che, nei test di laboratorio e nelle statistiche finanziarie, gli uomini mostrano in media una maggiore propensione al rischio. Ma quando la psicologa Bernd Figner e il suo team hanno analizzato più da vicino \textit{quali} rischi vengono presi e in \textit{quali} contesti, il quadro si è frantumato in un mosaico.\footnote{Figner, B., \& Weber, E. U. (2011). ``Who takes risks when and why? Determinants of risk taking''. \textit{Current Directions in Psychological Science}, 20(4), 211--216.} Gli uomini rischiano di più nel dominio finanziario e fisico (investimenti azzardati, sport estremi, guida spericolata). Ma nel dominio sociale (esporsi emotivamente, sfidare le norme del gruppo, affrontare conversazioni difficili), le donne rischiano almeno quanto, se non più degli uomini.

In altre parole, non esiste ``il sesso che rischia'' e ``il sesso che non rischia''. Esistono domini diversi in cui le persone, influenzate da biologia, cultura e contesto, calibrano il rischio in modo diverso. È come dire che i gatti non si arrampicano e i cani sì: vero in media per gli alberi, completamente falso se si parla di divani.

\section{L'empatia e il mito della lettura del pensiero}

Uno dei terreni più insidiosi del dibattito sulle differenze di genere riguarda l'empatia. Simon Baron-Cohen, neuroscienziato di Cambridge (e cugino dell'attore Sacha, per gli amanti delle curiosità), ha proposto una teoria influente: i cervelli maschili sarebbero, in media, più orientati alla ``sistematizzazione'' (comprendere regole, meccanismi, strutture), quelli femminili più orientati alla ``empatizzazione'' (comprendere emozioni, intenzioni, stati mentali).\footnote{Baron-Cohen, S. (2003). \textit{The Essential Difference: Men, Women, and the Extreme Male Brain}. Allen Lane.}

I dati a supporto sono reali ma modesti. Le donne ottengono punteggi mediamente più alti nei test di riconoscimento delle emozioni facciali. Ma la differenza, espressa in termini statistici, è piccola: la famosa \textit{d di Cohen} si aggira intorno a 0.2--0.5, il che significa una sovrapposizione enorme tra i due gruppi.\footnote{Thompson, A. E., \& Voyer, D. (2014). ``Sex differences in the ability to recognise non-verbal displays of emotion: A meta-analysis''. \textit{Cognition and Emotion}, 28(7), 1164--1195.} Per intenderci: prendete un uomo e una donna a caso dalla popolazione, e la probabilità che la donna sia più empatica dell'uomo è circa del 60\%, non del 100\%. Molto lontano dal determinismo biologico di ``Marte e Venere''.

Ma il problema non è tanto la differenza in sé. È ciò che la cultura ne fa. Quando una piccola tendenza statistica viene trasformata in destino biologico (``gli uomini non capiscono le emozioni'', ``le donne non capiscono la logica''), si innesca una profezia che si autoavvera. Il ragazzo che potrebbe sviluppare una fine sensibilità emotiva smette di provarci (``non è roba da maschi''). La ragazza che potrebbe eccellere in ingegneria si convince che non fa per lei (``non è roba da donne''). E le differenze, che all'inizio erano sfumature, si allargano fino a sembrare abissi.

È uno dei meccanismi più eleganti e più subdoli della psicologia sociale: l'effetto delle aspettative, o \textit{stereotype threat}.\footnote{Steele, C. M., \& Aronson, J. (1995). ``Stereotype threat and the intellectual test performance of African Americans''. \textit{Journal of Personality and Social Psychology}, 69(5), 797--811.} Claude Steele e Joshua Aronson lo dimostrarono per primi con le differenze razziali nei test cognitivi, ma il principio si applica con la stessa forza al genere: basta ricordare a qualcuno che ``il suo gruppo'' è considerato meno bravo in qualcosa, e le sue prestazioni peggioreranno. Non per incapacità, ma per il peso cognitivo dello stereotipo che si insinua nella mente come un rumore di fondo, consumando risorse mentali che dovrebbero andare al compito.

\section{La sindrome dell'impostore e il suo specchio}

Nel 1978, le psicologhe Pauline Clance e Suzanne Imes descrissero per la prima volta un fenomeno che osservavano nelle donne di successo: la convinzione persistente, nonostante le evidenze contrarie, di non meritare i propri risultati. Lo chiamarono \textit{impostor phenomenon}.\footnote{Clance, P. R., \& Imes, S. A. (1978). ``The imposter phenomenon in high achieving women: Dynamics and therapeutic intervention''. \textit{Psychotherapy: Theory, Research \& Practice}, 15(3), 241--247.}

Ricerche successive hanno mostrato che la sindrome colpisce anche gli uomini, ma con una frequenza e un'intensità significativamente inferiori. Il meccanismo è affascinante nella sua perversità: più una donna ha successo, più il suo cervello cerca spiegazioni alternative. Fortuna, tempismo, aiuto esterno, circostanze favorevoli. Qualsiasi cosa tranne la spiegazione più semplice: competenza.

Negli uomini, il pattern attributivo tende ad essere inverso. Il successo viene attribuito al talento personale, il fallimento alle circostanze esterne (sfortuna, contesto sfavorevole, colleghi incompetenti). Sono due bias speculari, due specchi deformanti: uno rimpicciolisce tutto, l'altro ingrandisce tutto.

Nessuno dei due riflette la realtà. Ma vivere con lo specchio che rimpicciolisce è particolarmente logorante, perché trasforma ogni successo in un'anomalia da spiegare, e ogni fallimento in una conferma di inadeguatezza. È come correre una maratona con uno zaino invisibile: arrivi lo stesso al traguardo, ma il costo energetico è enormemente superiore.

\section{Quando gli errori opposti si incontrano}

C'è un aspetto di questa storia che meriterebbe più attenzione di quanta ne riceva: cosa succede quando queste tendenze opposte si trovano nella stessa stanza?

La ricerca sulle dinamiche di gruppo offre una risposta illuminante. I team composti da persone con stili cognitivi diversi prendono, in media, decisioni migliori di quelli omogenei.\footnote{Woolley, A. W., et al. (2010). ``Evidence for a collective intelligence factor in the performance of human groups''. \textit{Science}, 330(6004), 686--688.} Ma, e questo è il punto cruciale, solo quando a tutti i membri viene data uguale possibilità di esprimersi. In assenza di questa condizione, prevale chi parla con più sicurezza, indipendentemente dalla qualità di ciò che dice. E statisticamente, nelle culture occidentali, chi parla con più sicurezza è più spesso un uomo.

Anita Woolley e il suo gruppo al MIT hanno scoperto qualcosa di ancor più specifico: il fattore che meglio prevede l'intelligenza collettiva di un gruppo non è il QI medio dei suoi membri, né quello del membro più brillante. È la ``sensibilità sociale media'' del gruppo, cioè la capacità dei suoi componenti di leggere gli stati emotivi degli altri e di distribuire equamente i turni di parola. E i gruppi con più donne? Tendevano ad avere punteggi di sensibilità sociale più alti e, di conseguenza, un'intelligenza collettiva superiore.

Non perché le donne siano ``più brave''. Ma perché un errore (eccesso di prudenza) si bilancia con l'altro (eccesso di audacia), e il risultato è una calibrazione migliore. A patto, naturalmente, che entrambe le voci vengano ascoltate.

\section{Il mosaico che cambia}

C'è una buona notizia in questa storia, e viene dai dati più recenti. I gap di genere nella fiducia, nella propensione al rischio, nella negoziazione si stanno restringendo.\footnote{Twenge, J. M. (2001). ``Changes in women's assertiveness in response to status and roles: A cross-temporal meta-analysis, 1931--1993''. \textit{Journal of Personality and Social Psychology}, 81(1), 133--145.} Le giovani generazioni mostrano differenze più contenute rispetto a quelle dei loro genitori. Non è che la biologia stia cambiando (l'evoluzione non lavora così in fretta): è la cultura che si sta spostando.

Questo dato, di per sé, ci dice qualcosa di fondamentale sulla natura di queste differenze. Se fossero interamente cablate nel cervello, non potrebbero cambiare nel giro di una generazione. Il fatto che cambino significa che una parte significativa di ciò che consideriamo ``naturale'' è in realtà il prodotto di aspettative sociali che, quando si modificano, modificano anche il comportamento.

Il cervello umano è un mosaico, come ci ha mostrato Daphna Joel. Non un mosaico statico, però: un mosaico che si ricompone continuamente in risposta all'ambiente. Le tessere biologiche esistono, ma il disegno che formano dipende in larga misura dal contesto in cui vengono posate.

La sfida, per ciascuno di noi, non è eliminare le differenze (sarebbe né possibile né auspicabile), ma riconoscere quando una differenza reale ma piccola viene amplificata dai nostri bias fino a diventare un muro. Quando il ``le donne tendono a'' diventa ``le donne sono''. Quando il ``gli uomini in media'' diventa ``tutti gli uomini sempre''. È in quello spazio, tra la sfumatura statistica e la caricatura culturale, che il danno si produce.

E forse la lezione più utile di questo capitolo non riguarda le differenze tra uomini e donne, ma qualcosa di più universale: il modo in cui il nostro cervello trasforma le sfumature in categorie nette, le probabilità in certezze, le tendenze in destini. È un bias che si applica al genere, ma anche all'età, alla nazionalità, alla professione, a qualsiasi gruppo umano. Il cervello ama le scorciatoie. La realtà, con testarda ostinazione, resta complicata.